\section{Provenance}\label{app:provenance}

The mathematics in \cref{sec:prelim,sec:densities,sec:identity,sec:pairs,sec:sawtooth,sec:single,sec:poly,sec:assembly} is an edited presentation of an archived run of a multi-agent proof pipeline in which language models proposed, proved and refereed the results. This appendix records where the material comes from. A proposition-by-proposition map, with every editorial change and correction, is in \texttt{manuscript/changes.md}.

\emph{Archive.}
\begin{itemize}
\item Source repository: \url{https://gitlab.com/mshepherd/squarefree-numbers}, archive label \texttt{zt-neutral}.
\item Archive commit: \texttt{ed0aebaf461b602601856bf1b693fd76893d9aba}.
\item Harness commit: \texttt{33aadc609359221ba5d81664438bce5f658f561c}.
\item Revision recorded for all run stages: \texttt{8dfc53d01e56c9949c42061fdbe907a98a12a5d7}.
\end{itemize}

\emph{Models and roles (as recorded).} The model identifiers are those recorded in the archive. The command-line tool versions are not model-weight versions, and model versions beyond these identifiers are unknown.
\begin{itemize}[before=\raggedright]
\item Planning and probing: \texttt{claude-opus-5-5} and \texttt{codex:gpt-6.1-sol}; probing also \texttt{claude-sonnet-5-5}.
\item Numerics: \texttt{claude-sonnet-5-5} and \texttt{codex:gpt-6.1-sol} produced candidate computations. The \texttt{codex:gpt-6.1-sol} snapshot was accepted.
\item Proofs: \texttt{claude-sonnet-5-5} (prover).
\item Refereeing: \texttt{claude-opus-5-5} and \texttt{codex:gpt-6.1-sol}. Each proposition was reviewed independently by both.
\item Integration (assembly audit): \texttt{claude-opus-5-5} and \texttt{codex:gpt-6.1-sol}, independently.
\end{itemize}
Claude models ran under Claude Code 2.1.309 and Codex models under codex-cli 0.166.0.

\emph{Write-up stage.} This text was written by \texttt{claude-opus-5-5} (Claude Code 2.1.311) from the archived run. It was refereed by \texttt{claude-opus-5-5} and \texttt{codex:gpt-6.1-sol} (the write-up stage records Claude Code 2.1.311 and codex-cli 0.166.0) in three rounds. Each round was followed by a revision, and the present version answers the third round. Some of the corrections described below were raised by these write-up referees.

\emph{Workflow.} The problem statement restricts $f$ to $\Z[t]$. It asks for the asymptotic behaviour of $V_f(x,y)$ in as large a range of $f(y)\le x$ as can be justified, and it accepts a weaker range or error term provided it is correct. The weight $f'(k)$ and the class~\eqref{eq:f-hyp} of polynomials used here are those of the problem statement; neither was changed. The numerical candidates were computed before any plan existed, and one was accepted on computational grounds only, not because of its fitted constants. Eight proof plans were written, and the plan \texttt{simplest@claude-opus-5-5} was selected by the project owner. It splits the argument into nine propositions P1--P9. In this text, P1 is \cref{prop:densities}, P2 is \cref{prop:identity}, P3 is \cref{prop:pairs}, P4 is \cref{prop:trivial}, P5 is \cref{prop:sawtooth}, P6 is \cref{prop:MT}, P7 is \cref{prop:single}, P8 is \cref{prop:poly}, and P9 is \cref{thm:main} with its proof in \cref{sec:assembly}. All nine proofs passed both referees in the first round, with only cosmetic remarks. Both integrators found that the propositions compose to the stated theorem without loss and without interface mismatches. The project owner's recorded decision reads: ``Let's just assume the proof is correct and complete.'' A separate read-only operator audit found no substantive gap and left final mathematical acceptance to the project owner. The corrections made in this text are local and fall into three groups, listed individually in \texttt{changes.md}: (a) those recorded in the operator audit (the normalisation and convergence wording in P1(iii), the endpoint $k=1$ in P8, and unused numerical remarks in P6 and P8); (b) cosmetic points raised by the run's referees; and (c) points raised at the write-up stage. Group~(c) contains one clarification of an archived auxiliary statement. For an infinite interval, the Euler--Maclaurin lemma in the archived P5 proof says that ``the same holds'' under the hypotheses $G,G'\to0$, $G''\in L^1$. Read as an identity between a convergent series and a convergent integral, this is false in that generality. \Cref{lem:EM} keeps exactly the archived hypotheses and states what letting $b\to\infty$ gives: the difference between the partial sums and the partial integrals converges, and the formula holds whenever the series or the integral converges. Both archived applications have absolutely convergent series and integrals, so no proposition and no step of the argument changes. No decision of the project owner on this point is recorded, and \texttt{changes.md} lists it for the owner's attention. The other items in group~(c) concern only the exposition of this text, for example an arithmetic slip in an earlier version of the write-up, and do not touch the archived arguments. Agreement between models is evidence, not certification; the arguments are reproduced here in full so that they can be checked directly.

\emph{Independent checks during the run.} Before acceptance, the operator compared both numerical candidates on overlapping rows and recomputed some rows by brute force. One integrator recomputed all $16$ rows for $k^2+1$ and $k^3+2$ at $x=250000$ by direct counting. The other computed $\sum f'(k)\MT(x;f(k))$ independently and compared it with the data. These checks are recorded in the archive. The computations in \cref{sec:numerics} were redone for this write-up as described there.

\emph{Literature.} The run used the papers supplied with the problem and literature up to 2010, with the single exception of Brüdern--Wooley~\cite{BruedernWooley2011}, which was supplied. The present text cites only works in the supplied corpus. See \texttt{README.md} for the limitations of the bibliography.
